\begin{abstract}
A thermal convection loop is a circular chamber filled with water, heated on the bottom half and cooled on the top half.
With sufficiently great forcing of heat, the direction of fluid flow in the loop oscillates chaotically, thus forming an analog to the real weather system.
Like full scale weather models, we can only observe a small part of state space, making prediction very difficult.
To overcome this challenge, data assimilation methods and specifically ensemble methods use the computational model itself to estimate the uncertainty of our model to optimally combine these observations into an initial condition for predicting the future.
First, we build and verify four distinct DA methods.
Then, a computational fluid dynamic simulation of the loop and a reduced order model are both used by these DA methods to predict flow reversals.
This happens in real time as experimental data is fed into the algorithm.

\end{abstract}
